\section{问题三：预报融合与多阶段滚动调整}
\subsection{发布预报的处理和在线融合}
附件3提供每次发布后24小时的整点光伏预报。先线性插值至十分钟网格，发布瞬间端点采用当时已观测光伏，0点采用上一日末观测。超过预报覆盖范围的次日时段用历史预测补足。负载仍由历史模型预测，不以当天尚未发生的实际负载更新。

设同一发布时刻、同一未来六小时块的历史预测为$H$、发布预报为$F$、实际光伏为$V$。采用最近56天内已经完整验证的预报样本，计算
\begin{equation}
\alpha=\operatorname{clip}_{[0,1]}
\frac{\sum(F-H)(V-H)}{\sum(F-H)^2}.\label{eq:fusion}
\end{equation}
以$\alpha F+(1-\alpha)H$为基础，再加历史平均残差乘$n/(n+14)$的收缩修正，并裁剪至非负。分母退化时按原实现回退；有效样本少于7条时暂用发布预报。每个发布时刻与未来六小时块分别估计，避免把夜间和日间的误差结构混为一谈。

跨午夜预报行只有在未来24小时全部实现后才能训练。因此，决策日使用的融合训练行最晚发布于日前两天；当天后续预报不能反向参与当前权重估计。56天窗口和收缩常数14沿用既有参数，不声称经过独立测试选型。

\fig{05_q3_forecast_fusion}
图\ref{fig:05_q3_forecast_fusion}显示，6—12时段的MAE由历史预测309.24 kW降至融合预测169.41 kW，12—18时段由288.99 kW降至133.68 kW。日夜面板使用不同纵轴尺度。融合包含偏差修正，不能仅将其理解为两条曲线的固定凸组合；预测误差改善也不直接等于同比例费用下降。

\subsection{场景树与非前视决策}
采用与式\eqref{eq:residual}同类的整日残差场景，但点预测和历史残差都对应当前发布时刻。每个历史日还携带后续预报修订向量，用确定性二叉聚类近似下一次发布可能带来的信息，最多形成四层、15个决策节点。聚类依据是届时可见的预报修订，不是尚未观测的实际负载或实际光伏。

设节点$n$包含历史场景集合$\mathcal S_n$，概率为$\pi_n=|\mathcal S_n|/S$。同一节点只设置一组购电、充放电和库存变量，各场景仅缺口补救不同。因此，无法用当前信息区分的场景自动共享决策，满足非前视要求。子节点的初始储电量等于父节点末值；仍让全部历史样本参与缺口期望，不以少量均值曲线替代全部误差样本。

\subsection{调整合同与优化模型}
设$g_t^0$为0点原计划，$q_t$为最终执行购电，
\begin{equation}
u_t=\max(q_t-g_t^0,0),\qquad v_t=\max(g_t^0-q_t,0).
\end{equation}
采用取消退原费并另付50\%违约费的解释，则实际费用为
\begin{align}
J_3&=\sum_t\left(p_tg_t^0+1.5p_tu_t-p_tv_t+0.5p_tv_t+5p_te_t\right)\notag\\
&=\sum_t\left(p_tg_t^0+1.5p_tu_t-0.5p_tv_t+5p_te_t\right).\label{eq:q3settle}
\end{align}
退款不是负购电，净调整费可以为负。原计划费、净调整费和紧急费相加时，不再重复计入退款、违约费与增购费。题面使用“交易时刻电价”，现有模型按购电所属供电时段定价；如改为决策发布时刻定价，需改变优化与结算，不能只调整文字。

0点原计划跨场景共享，节点变量满足$q_{n,t}=g_t^0+u_{n,t}-v_{n,t}$以及
\begin{equation}
q_{n,t}+b_{n,t}-c_{n,t}+e_{s,t}\ge N_{s,t},\qquad s\in\mathcal S_n.
\end{equation}
另保留净计划供给不低于节点点预测的约束和共同储能条件。省略极小的吞吐量消歧项，0点优化目标为
\begin{equation}
\min\ \sum_t p_tg_t^0+
\sum_n\pi_n\sum_{t\in\mathcal T_n}(1.5p_tu_{n,t}-0.5p_tv_{n,t})
+\frac5S\sum_s\sum_t p_te_{s,t}+\Phi.\label{eq:q3obj}
\end{equation}
其中$\mathcal T_n$为节点执行区间，$\Phi$为各叶节点接入的次日确定性前瞻费用。0点第一执行块不调整，后续更新时原计划固定，原计划费作为常数从优化中省略，但实际结算仍完整保留。

普通日的次日末端固定为本日0点储电量，6、12、18点重算时不随当前库存改变该终端目标。12月31日直接固定当日末6000 kWh。每次仅执行本发布时刻至下一允许更新时刻的控制，其后计划只是预案；每个供电时段最多实际调整一次，不依赖反复买卖后仅结算净差额的假设。

\subsection{更新频率的经济价值}
\begin{table}[htbp]\centering\small
\caption{第三问更新频率对照}
\begin{tabular}{lrr}\toprule
方案 & 总费用/元 & 紧急购电量/kWh\\\midrule
同口径历史预测 & 14389892.50 & 246099.33\\
仅0点融合 & 14181445.03 & 218589.30\\
0、6点 & 14101371.47 & 186570.89\\
0、6、12点 & 14022358.10 & 175344.39\\
四次更新主方案 & 14022396.98 & 175369.95\\
\bottomrule\end{tabular}\end{table}

同口径历史预测对照为14389892.50元，其初始化和物理处理与本问一致，并非第二问正式结果14389135.08元。四次更新主方案总费用14022396.98元，节省367495.52元，紧急购电量由246099.33 kWh降至175369.95 kWh。

\fig{04_q3_information_value}
图\ref{fig:04_q3_information_value}把改善分为引入0点融合和后续新增调整。仅0点融合节省208447.47元；增加6点、12点分别再节省80073.56元和79013.36元；加入18点后费用反增38.88元。故本年度数据支持重视6点和12点调整，未显示18点的明确额外经济收益，预设四次更新主结果仍予保留。

这不意味着18点信息始终无价值。它可能改变次日预报与库存安排，且本方法在有限历史样本上重建树并滚动执行，真实费用并不保证随更新机会增加而下降。不能利用同年最好的更新组合，反向宣称它是事先已知的最优规则。

\begin{table}[htbp]\centering\small
\caption{第三问主结果实际费用分解}
\begin{tabular}{lr}\toprule
项目 & 费用/元\\\midrule
原计划费 & 13345819.30\\
增购费 & 89947.28\\
取消退款（扣减） & -150170.13\\
违约费 & 75085.07\\
净调整费 & 14862.22\\
紧急费 & 661715.46\\
总费用 & 14022396.98\\
\bottomrule\end{tabular}\end{table}

\subsection{指定日期的日内行为}
\fig{06_q3_rolling_adjustment}
图\ref{fig:06_q3_rolling_adjustment}展示四个指定日期最终执行购电减0点原计划的功率差，并对齐储电量和紧急购电时段。正值表示增购，负值表示取消；橙色短线只标记紧急购电发生时间，不表示数量。该图不是全部中间计划的修订历史。库存轨迹说明，日内调整既影响当前缺口，也受到后续用电和终端条件制约。
\FloatBarrier
